Llano es una skill para agentes de código que adapta ASD-STE100, el inglés técnico simplificado de la aviación, al español. Este piloto compara llano con una instrucción simple de concisión, con respuestas sin instrucción y con caveman. Usa \evalPrompts{} prompts técnicos en español y Claude \evalModel{} dentro de Claude Code.

Llano no acorta las respuestas frente a la instrucción de concisión: la diferencia en palabras es \deltaWordsTerse{} (IC 95\%: \deltaWordsTerseLo{} a \deltaWordsTerseHi{}). Su efecto está en la forma. Baja las violaciones de reglas de \tersePerHundred{} a \llanoPerHundred{} cada 100 palabras, con menos violaciones en \violBetterTerse{} de \evalPrompts{} prompts ($p$ \violPTerse). También sube la legibilidad INFLESZ de \terseInflesz{} a \llanoInflesz. La fidelidad no cambia: llano conserva el \llanoFidelity{} de los hechos clave, igual que la instrucción de concisión. En la comparación a ciegas, el juez prefirió llano en \winsTerse{} de \evalPrompts{} casos, pero la diferencia no es significativa ($p = \pSignTerse$). El costo es la entrada: cada llamada suma unos \llanoInputTokens{} tokens contra \terseInputTokens.
